\section{Programa principal (host)}
\begin{lstlisting}[title={\textit{Macros del programa}}, captionpos=b, language=c, firstnumber=9]
#define N_REFERENCIA 512
#define DISTANCIA_REFERENCIA 20
#define TEMP_FUENTE 100.0f
#define DIFUSION 0.25f
#define ANCHO_BLOQUE 32
\end{lstlisting}

\subsection{Validación de parámetros}
\begin{lstlisting}[language=c, firstnumber=15]
bool validar_parametros(int argc, char *argv[])
\end{lstlisting}
Esta función se asegura de que:
\begin{itemize}
  \item se pasen los tres parámetros <N> <Pasos de simulación> <hilos por bloque>;
  \item N sea potencia de 2 \footnote{Requerimiento del enunciado.};
  \item N sea mayor o igual a 32 \footnote{Para que las fuentes de calor respeten la proporcionalidad del algoritmo secuencial ya dado sin estar tan juntas entre sí.};
  \item se ingrese una cantidad de pasos mayor o igual a 1;
  \item los hilos por bloque sean múltiplos de 32, entre 32 y 1024 \footnote{Las GPUs NVIDIA, desde compute capability 2.0, aceptan un máximo de 1024 hilos por bloque.}.
\end{itemize}
\label{sec:explicacion_ancho_bloque}
\subsection{Preparación de la grilla}
    \subsubsection{Distancia entre las fuentes de calor}
    En el algoritmo secuencial, existe una fuente de calor en el centro de la grilla y 4 más en diagonal, a una distancia fija de 20 celdas en $x$ y en $y$.\\
    En esta implementación paralela decidí que esa distancia no sea fija ya que, en una grilla grande, las fuentes quedan muy juntas y, en una grilla muy chica, no entran.
    De esta forma, se sigue respetando la proporción de 512 celdas $\longleftrightarrow$ distancia 20 del programa secuencial:
\begin{lstlisting}[language=c, firstnumber=161]
int distancia = N * DISTANCIA_REFERENCIA / N_REFERENCIA;
\end{lstlisting}
    
    \subsubsection{Gestión inicial de memoria}
    \begin{enumerate}
        \item Calcular la cantidad total de bytes que ocupará la grilla:
\begin{lstlisting}[language=c, firstnumber=163]
size_t bytes = (size_t)N * N * sizeof(float);
\end{lstlisting}
        \item Reservar memoria en el host para la grilla:
\begin{lstlisting}[language=c, firstnumber=164]
float *h_grilla = (float *)malloc(bytes);
\end{lstlisting}
\textit{Una grilla de dos dimensiones se guarda en un arreglo de una dimensión, fila tras fila (ordenamiento por filas).}
        \item Reservar memoria en GPU para las grillas:
\begin{lstlisting}[language=c, firstnumber=166]
float *d_grilla;
HANDLE_ERROR(cudaMalloc(&d_grilla, bytes));
float *d_nueva;
HANDLE_ERROR(cudaMalloc(&d_nueva, bytes));
\end{lstlisting}
        \item Utilizo cudaMemset para evitar la transferencia de la grilla del host hacia las grillas en GPU, haciendo que esta las inicialice directamente:
\begin{lstlisting}[language=c, firstnumber=181]
HANDLE_ERROR(cudaMemset(d_grilla, 0, bytes));
HANDLE_ERROR(cudaMemset(d_nueva, 0, bytes));
\end{lstlisting}
    \end{enumerate}

\subsubsection{Fuentes de calor iniciales}
La inicialización de las 5 fuentes de calor en 100 la lleva a cabo un kernel (directamente la GPU), y el host solo recibe copias para guardar los frames:

\begin{lstlisting}[language=c, firstnumber=49]
__global__ void inicializar_fuentes_de_calor(float *grilla, int N, int distancia) {
    int centroX = N / 2;
    int centroY = N / 2;

    grilla[centroY * N + centroX] = TEMP_FUENTE;
    grilla[(centroY + distancia) * N + (centroX + distancia)] = TEMP_FUENTE;
    grilla[(centroY - distancia) * N + (centroX + distancia)] = TEMP_FUENTE;
    grilla[(centroY + distancia) * N + (centroX - distancia)] = TEMP_FUENTE;
    grilla[(centroY - distancia) * N + (centroX - distancia)] = TEMP_FUENTE;
}
\end{lstlisting}

\subsubsection{Dimensionamiento de hilos y bloques}
\begin{itemize}
    \item \textbf{Dimensiones de bloque:} este ancho de bloque está fijado en 32 hilos para que coincida con el tamaño del warp. De esta forma, cada fila del bloque corresponde exactamente a un warp, donde sus hilos contiguos ocupan celdas contiguas de dicha fila. Al ser 32 celdas de 4 bytes (float) por fila, cada warp accede a 128 bytes contiguos que coinciden exactamente con el tamaño de segmento de memoria global para palabras de 4 bytes. \\
    Por otro lado, al ser N potencia de 2 y mayor o igual a 32, los accesos del warp quedan alineados con cada segmento (o sea, que empiezan en una dirección múltiplo de 128). Así, cada lectura o escritura de los bloques se resuelve de forma coalescente.
    
    Este mismo ancho también sirve para el acceso a memoria compartida, que se divide en 32 bancos de palabras de 4 bytes c/u. Como los 32 hilos de cada warp acceden a 32 floats consecutivos de una misma fila, cada uno accede a un banco distinto y el acceso se resuelve sin necesidad de serializar (lo cual costaría hasta 32 accesos por warp).
\begin{table}[H]
\centering
    \begin{tabular}{rrr}
        \toprule
                     & \multicolumn{2}{c}{Bloque} \\
        \cmidrule(lr){2-3}
        hilosXBloque & Ancho & Alto \\
        \midrule
        32   & 32 & $\frac{32}{32} = 1$ \\
        \addlinespace
        64   & 32 & $\frac{64}{32} = 2$ \\
        \addlinespace
        128  & 32 & $\frac{128}{32} = 4$ \\
        \addlinespace
        256  & 32 & $\frac{256}{32} = 8$ \\
        \addlinespace
        512  & 32 & $\frac{512}{32} = 16$ \\
        \addlinespace
        1024 & 32 & $\frac{1024}{32} = 32$ \\
        \bottomrule
    \end{tabular}
\end{table}
\begin{lstlisting}[language=c, firstnumber=192]
dim3 bloque(ANCHO_BLOQUE, hilosXBloque / ANCHO_BLOQUE);
\end{lstlisting}
    
    \item \textbf{Dimensiones de grid:} hacen falta suficientes bloques para cubrir las N columnas y filas. En algunos casos las filas pueden sobrar, y esos son los hilos que no van a operar sobre ninguna celda de la grilla:
\begin{table}[H]
  \centering
  \small
  \begin{tabular}{rrrrrrl}
    \toprule
      &              & \multicolumn{2}{c}{Bloque} & \multicolumn{2}{c}{Grid} & \\
    \cmidrule(lr){3-4} \cmidrule(lr){5-6}
    N & hilosXBloque & Ancho & Alto & Ancho & Alto & Hilos totales \\
    \midrule
    32   & 256  & 32 & $\frac{256}{32} = 8$   & $\left\lceil \frac{32}{32} \right\rceil = 1$    & $\left\lceil \frac{32}{8} \right\rceil = 4$      & $32 \times 32$ (justo) \\
    \addlinespace
         & 96   & 32 & $\frac{96}{32} = 3$    & $\left\lceil \frac{32}{32} \right\rceil = 1$    & $\left\lceil \frac{32}{3} \right\rceil = 11$     & $32 \times \mathbf{33}$ (sobra 1 fila) \\
    \addlinespace
    64   & 128  & 32 & $\frac{128}{32} = 4$   & $\left\lceil \frac{64}{32} \right\rceil = 2$    & $\left\lceil \frac{64}{4} \right\rceil = 16$     & $64 \times 64$ (justo) \\
    \addlinespace
         & 160  & 32 & $\frac{160}{32} = 5$   & $\left\lceil \frac{64}{32} \right\rceil = 2$    & $\left\lceil \frac{64}{5} \right\rceil = 13$     & $64 \times \mathbf{65}$ (sobra 1 fila) \\
    \addlinespace
    128  & 512  & 32 & $\frac{512}{32} = 16$  & $\left\lceil \frac{128}{32} \right\rceil = 4$   & $\left\lceil \frac{128}{16} \right\rceil = 8$    & $128 \times 128$ (justo) \\
    \addlinespace
         & 192  & 32 & $\frac{192}{32} = 6$   & $\left\lceil \frac{128}{32} \right\rceil = 4$   & $\left\lceil \frac{128}{6} \right\rceil = 22$    & $128 \times \mathbf{132}$ (sobran 4 filas) \\
    \addlinespace
    256  & 1024 & 32 & $\frac{1024}{32} = 32$ & $\left\lceil \frac{256}{32} \right\rceil = 8$   & $\left\lceil \frac{256}{32} \right\rceil = 8$    & $256 \times 256$ (justo) \\
    \addlinespace
         & 224  & 32 & $\frac{224}{32} = 7$   & $\left\lceil \frac{256}{32} \right\rceil = 8$   & $\left\lceil \frac{256}{7} \right\rceil = 37$    & $256 \times \mathbf{259}$ (sobran 3 filas) \\
    \addlinespace
    512  & 256  & 32 & $\frac{256}{32} = 8$   & $\left\lceil \frac{512}{32} \right\rceil = 16$  & $\left\lceil \frac{512}{8} \right\rceil = 64$    & $512 \times 512$ (justo) \\
    \addlinespace
         & 96   & 32 & $\frac{96}{32} = 3$    & $\left\lceil \frac{512}{32} \right\rceil = 16$  & $\left\lceil \frac{512}{3} \right\rceil = 171$   & $512 \times \mathbf{513}$ (sobra 1 fila) \\
    \addlinespace
    1024 & 512  & 32 & $\frac{512}{32} = 16$  & $\left\lceil \frac{1024}{32} \right\rceil = 32$ & $\left\lceil \frac{1024}{16} \right\rceil = 64$  & $1024 \times 1024$ (justo) \\
    \addlinespace
         & 640  & 32 & $\frac{640}{32} = 20$  & $\left\lceil \frac{1024}{32} \right\rceil = 32$ & $\left\lceil \frac{1024}{20} \right\rceil = 52$  & $1024 \times \mathbf{1040}$ (sobran 16 filas) \\
    \bottomrule
  \end{tabular}
  \label{tab:grid}
\end{table}
\begin{lstlisting}[language=c, firstnumber=193]
dim3 grid((N + bloque.x - 1) / bloque.x, (N + bloque.y - 1) / bloque.y);
\end{lstlisting}
\end{itemize}

\subsection{Bucle de pasos}
En cada paso primero se guarda en disco el estado actual de la grilla y después se lanza el kernel, que calcula el paso siguiente en d\_nueva. Para que ese resultado sea la entrada del próximo paso no copio nada, sino que intercambio los punteros, así d\_grilla pasa a apuntar a la grilla recién calculada y d\_nueva a la vieja, que se va a sobrescribir.

Por último copio la grilla al host para poder guardarla en la siguiente vuelta. En el último paso esa copia no se hace, porque ya no queda ningún frame por guardar. Las líneas que omití en el fragmento son las de medición de tiempos.

\begin{lstlisting}[firstnumber=196]
for (int paso = 0; paso < pasos; paso++) {(*@\setcounter{lstnumber}{197}@*)
    guardar_paso(h_grilla, paso, N);
    ...(*@\setcounter{lstnumber}{201}@*)
    difusion<<<grid, bloque, bytes_compartidos>>>(d_grilla, d_nueva, N, distancia);
    ...(*@\setcounter{lstnumber}{208}@*)
    float *aux = d_grilla;
    d_grilla = d_nueva;
    d_nueva = aux;

    if (paso + 1 < pasos) {
        ...
        HANDLE_ERROR(cudaMemcpy(h_grilla, d_grilla, bytes, cudaMemcpyDeviceToHost));
        ...
    }
}
\end{lstlisting}